\papersection{Introduction}{sec:intro}

Scheduled macroeconomic releases are among the cleanest information events in financial markets.
Their timing is announced months in advance, and futures prices react within the first minute
\citep{Fleming1999, Balduzzi2001, Andersen2003}. In equity-index futures, however, part of the
move continues well after that minute: after a hot CPI print the E-mini S\&P 500 (\ES) falls about
37 basis points (bps) in the first minute and a further 17 bps over the next half hour. This paper
asks whether the \emph{shape} of the market's reaction to a release, compared with its reactions
to earlier releases, forecasts that next move.

We match the price path around each release to past paths with dynamic time warping (DTW) and
trade the move that followed the closest matches. The book trades five back-to-back 30-minute
windows after every release that moves the market, enters only after the release is public, and
is scored on the held-out years 2021--2026 after costs in four futures markets: the S\&P 500 (\ES),
Nasdaq-100 (\NQ) and Dow (\YM) equity indices and the 10-year Treasury note (\ZN).

Three findings stand out. First, the dispersion of the matches forecasts the size of the next
move in every market and year, while their direction is weak and changed sign around 2018.
Second, the \ES\ book earns a Sharpe ratio of 0.64 at a quarter tick per side with no market beta,
but \NQ\ earns only 0.11 and \YM, added later as an out-of-sample market, loses ($-0.39$). Third,
the edge is concentrated in the first window after the release, and in \ZN\ it is smaller than
one tick. The Sensibility part explains the approach, the Data and Model parts describe the inputs
and the strategy, and the Backtest part reports the results.
